\chapterimage{chapter_head_2.pdf} % Chapter heading image
\label{Apendice B}

\chapter{Ecuaciones de Campo E-L}

En este apéndice se da a conocer una manera alternativa de derivar las ecuaciones de Euler-Lagrange.
Se considera un conjunto uniparam\'{e}trico de funciones variadas que se reducen al valor correcto de $\phi
_{r}\left( x \right) $ cuando $\alpha $ va a cero. Una familia de estas funciones se puede constru\'{\i}r como:

\begin{equation*}
\phi_{r}\left( x,\alpha \right) =\phi_{r}\left( x,0\right) +\alpha \zeta _{r}\left( x\right) ~,
\end{equation*}

las funciones $\zeta_{r}\left(x\right)$ son funciones de clase $C^{2}$ en todos los par\'{a}metros. Dichas funciones se deben anular sobre la hipersuperficie $\Sigma$ que delimita a la regi\'{o}n de integraci\'{o}n. La anulaci\'{o}n del variacional de $S$, es equivalente a establecer
que la derivada de $S$ con respecto a $\alpha $ sea cero.

\begin{equation}
\frac{dS}{d\alpha }=\frac{d}{d\alpha }\int d^{4}x\ \mathcal{L}\left(\phi_{r},\partial_{\mu }\phi_{r},x\right) =\int d^{4}x\
\left\{ \frac{\partial \mathcal{L}}{\partial \phi_{r}}\frac{\partial \phi _{r}}{\partial \alpha }\mathcal{+}\frac{\partial \mathcal{L}}{\partial \left( \partial _{\mu }\phi_{r}\right) }\frac{\partial \left(\partial _{\mu }\phi _{r}\right) }{\partial \alpha }\right\} ~.
\label{2.19}
\end{equation}
integrando por partes se obtiene,
\begin{equation}
\frac{dS}{d\alpha }=\int d^{4}x\ \frac{\partial \phi_{r}}{\partial \alpha }\left\{ \frac{\partial \mathcal{L}}{\partial \phi_{r}}-\partial
_{\mu}\left( \frac{\partial \mathcal{L}}{\partial \left( \partial _{\mu }\phi_{r}\right) }\right) \right\} +\int d^{4}x\partial_{\mu}\left(\frac{\partial \mathcal{L}}{\partial \left( \partial _{\mu }\phi_{r}\right) }\frac{\partial \phi_{r}}{\partial \alpha }\right)
~.  \label{2.19a}
\end{equation}

La segunda integral se elimina cuando $\alpha $ tiende a cero, lo cual se
puede ver de varias formas:
\vspace{0.5cm}

\begin{itemize}
\item Si se observa t\'{e}rmino a t\'{e}rmino, al realizar la integraci\'{o}n sobre un parámetro $x$ en particular para cada t\'{e}rmino de la derivada, se v\'{e} que estos t\'{e}rminos se tienen que desaparecer porque la derivada con respecto a $\alpha $ de los campos es cero en los extremos.\\


\item La integral se puede transformar por medio de un teorema de la
divergencia en 4 dimensiones en una integral sobre la hipersuperficie que
delimita la regi\'{o}n de integraci\'{o}n en el cuadriespacio (\textbf{euclidiano})\footnote{Procedimiento realizado en el desarrollo teórico en la variación de la acción}. La integral de superficie tambi\'{e}n se anula porque la variaci\'{o}n de $\phi_{r}$ en la vecindad de la funci\'{o}n de campo correcta es cero sobre la superficie. Para ver esto se calcula la segunda integral a la derecha en (\ref{2.19a}), as\'{\i}
\end{itemize}

\begin{eqnarray}
\int d^{4}x~\partial _{\mu }\left( \frac{\partial \mathcal{L}}{\partial
\left( \partial _{\mu }\phi _{r}\right) }\frac{\partial \phi _{r}}{\partial
\alpha }\right)  &=&\int d^{4}x~\left[ \partial _{0}\left( \frac{\partial 
\mathcal{L}}{\partial \left( \partial _{0}\phi _{r}\right) }\frac{\partial
\phi _{r}}{\partial \alpha }\right) +\partial _{i}\left( \frac{\partial 
\mathcal{L}}{\partial \left( \partial _{i}\phi _{r}\right) }\frac{\partial
\phi _{r}}{\partial \alpha }\right) \right]   \notag \\
&=&\int d^{4}x\partial _{0}\left( \frac{\partial \mathcal{L}}{\partial
\left( \partial _{0}\phi _{r}\right) }\frac{\partial \phi _{r}}{\partial
\alpha }\right) +\int d^{4}x\partial _{i}\left( \frac{\partial \mathcal{L}}{%
\partial \left( \partial _{i}\phi _{r}\right) }\frac{\partial \phi _{r}}{%
\partial \alpha }\right)   \notag \\
&=&\int d^{3}x\int_{t_{1}}^{t_{2}}dt~\partial _{0}\left( \frac{\partial 
\mathcal{L}}{\partial \left( \partial _{\mu }\phi _{r}\right) }\frac{%
\partial \phi _{r}}{\partial \alpha }\right)   \notag \\
&&+\int_{t_{1}}^{t_{2}}dt~\int d^{3}x~\partial _{i}\left( \frac{\partial 
\mathcal{L}}{\partial \left( \partial _{i}\phi _{r}\right) }\frac{\partial
\phi _{r}}{\partial \alpha }\right)   \notag \\
&=&\int d^{3}x~\left. \left( \frac{\partial \mathcal{L}}{\partial \left(
\partial _{\mu }\phi _{r}\right) }\frac{\partial \phi _{r}}{\partial \alpha }%
\right) \right\vert _{t_{1}}^{t_{2}}  \notag \\
&&+\int_{t_{1}}^{t_{2}}dt~\int d^{3}x~\partial _{i}\left( \frac{\partial 
\mathcal{L}}{\partial \left( \partial _{i}\phi _{r}\right) }\frac{\partial
\phi _{r}}{\partial \alpha }\right) ~.
\end{eqnarray}

De este resultado se establece que la primera integral se anula por la
condici\'{o}n de frontera $\left. \frac{\partial \phi_{r}}{\partial
\alpha }\right\vert _{t=t_{1}}=\left. \frac{\partial \phi_{r}}{\partial
\alpha }\right\vert _{t=t_{2}}=0$. De igual manera, en la segunda integral
el integrando esta siendo evaluado en la superficie que encierra el volumen en el cual est\'{a} siendo considerados los campos. Si este volumen se extiende en todo el espacio, la superficie que lo limita tiende a infinito y del comportamiento asint\'{o}tico de los campos hace que esta integral de superficie tienda a cero, es decir

\begin{equation*}
\int_{t_{1}}^{t_{2}}dt\int d^{3}x~\partial _{i}\left( \frac{\partial \mathcal{L}}{\partial \left( \partial _{i}\phi_{r}\right) }\frac{\partial\phi_{r}}{\partial \alpha }\right)=\int_{t_{1}}^{t_{2}}dt\oint_{S}ds_{i}\left(\frac{\partial \mathcal{L}}{\partial \left( \partial _{i}\phi_{r}\right) }\zeta_{r}\right) =0~\  .
\end{equation*}

\noindent
De esta manera, se deduce que la variaci\'{o}n 

\begin{equation}
\delta \int d^{4}x\ \partial _{\mu }\left( \frac{\partial \mathcal{L}}{%
\partial \left( \partial _{\mu }\phi_{r}\right) }\frac{\partial \phi
_{r}}{\partial \alpha }\right) =\delta \int d^{4}x\ \partial _{\mu
}\left( \frac{\partial \mathcal{L}}{\partial \left( \partial _{\mu }\phi
_{r}\right) }\zeta _{r}\right) =0 ~.  \label{2.21}
\end{equation}

\noindent
Por lo anterior, en el l\'{\i}mite cuando $\alpha \rightarrow 0$, la ecuaci\'{o}n (\ref{2.19a}) se reduce a,

\begin{equation}
\left( \frac{dS}{d\alpha }\right) _{\alpha =0}=\int d^{4}x\ \left( \frac{\partial \phi_{r}}{\partial \alpha }\right) _{\alpha =0}\left\{ \frac{\partial \mathcal{L}}{\partial \phi _{r}}-\partial _{\mu }\left( \frac{\partial \mathcal{L}}{\partial \left( \partial _{\mu }\phi_{r}\right) }\right) \right\} =0~,  \label{2.21a}
\end{equation}
y de la naturaleza arbitraria de la variaci\'{o}n de cada campo $\phi_{r}$, la ecuaci\'{o}n (\ref{2.21a}) se satisface s\'{o}lo si se anula el t\'{e}rmino entre llaves, entonces del lema fundamental del c\'{a}lculo de variaciones se determina,

\begin{equation}
\frac{\partial \mathcal{L}}{\partial \phi_{r}}-\partial _{\mu }\left(\frac{\partial \mathcal{L}}{\partial \left( \partial _{\mu }\phi _{r}\right) }\right) =0~.  \label{2.22}
\end{equation}

La relaci\'{o}n (\ref{2.22}) es conocida como ecuaci\'{o} de
Euler-Lagrange, y representa un conjunto de ecuaciones diferenciales
parciales para las variables de campo, el n\'{u}mero de ecuaciones est\'{a} determinado por los valores posibles de $r$. Vale la pena enfatizar de nuevo que dado que cada $x$ corresponde a un \'{\i}ndice continuo para las variables de campo. Cada una de las ecuaciones (\ref{2.22}) para $r$ fijo, corresponde a un conjunto completo de ecuaciones diferenciales de movimiento de Lagrange en el caso discreto.

\begin{exercise}
Utilizar integraci\'{o}n por partes para demostrar la ecuaci\'{o}n (\ref%
{2.19a}).
\end{exercise}

\begin{solution}
Al utilizar la integraci\'{o}n por partes%
\begin{eqnarray*}
\int dx^{\mu }\partial _{\nu }\left( \frac{\partial \mathcal{L}}{\partial
\eta _{\rho ,\nu }}\frac{\partial \eta _{\rho }}{\partial \alpha }\right)
&=&\int dx^{\mu }\partial _{\nu }\left( \frac{\partial \mathcal{L}}{\partial
\left( \partial _{\nu }\eta _{\rho }\right) }\right) \frac{\partial \eta
_{\rho }}{\partial \alpha }+\int dx^{\mu }\frac{\partial \mathcal{L}}{%
\partial \left( \partial _{\nu }\eta _{\rho }\right) }\partial _{\nu }\left( 
\frac{\partial \eta _{\rho }}{\partial \alpha }\right) \\
&=&\int dx^{\mu }\partial _{\nu }\left( \frac{\partial \mathcal{L}}{\partial
\left( \partial _{\nu }\eta _{\rho }\right) }\right) \frac{\partial \eta
_{\rho }}{\partial \alpha }+\int dx^{\mu }\frac{\partial \mathcal{L}}{%
\partial \eta \left( \partial _{\nu }\eta _{\rho }\right) }\frac{\partial
\left( \partial _{\nu }\eta _{\rho }\right) }{\partial \alpha }~~~~~,
\end{eqnarray*}%
de manera que:%
\begin{equation*}
\int dx^{\mu }\frac{\partial \mathcal{L}}{\partial \left( \partial _{\nu
}\eta _{\rho }\right) }\frac{\partial \left( \partial _{\nu }\eta _{\rho
}\right) }{\partial \alpha }=\int dx^{\mu }\left[ \partial _{\nu }\left( 
\frac{\partial \mathcal{L}}{\partial \left( \partial _{\nu }\eta _{\rho
}\right) }\frac{\partial \eta _{\rho }}{\partial \alpha }\right) -\partial
_{\nu }\left( \frac{\partial \mathcal{L}}{\partial \left( \partial _{\nu
}\eta _{\rho }\right) }\right) \frac{\partial \eta _{\rho }}{\partial \alpha 
}\right] ~~~~~,
\end{equation*}%
lo que implica en (\ref{2.19})%
\begin{eqnarray}
\frac{dS}{d\alpha } &=&\int dx^{\mu }\ \left\{ \frac{\partial \mathcal{L}}{%
\partial \eta _{\rho }}\frac{\partial \eta _{\rho }}{\partial \alpha }%
\mathcal{+}\partial _{\nu }\left( \frac{\partial \mathcal{L}}{\partial
\left( \partial _{\nu }\eta _{\rho }\right) }\frac{\partial \eta _{\rho }}{%
\partial \alpha }\right) -\partial _{\nu }\left( \frac{\partial \mathcal{L}}{%
\partial \left( \partial _{\nu }\eta _{\rho }\right) }\right) \frac{\partial
\eta _{\rho }}{\partial \alpha }\right\}  \notag \\
&=&\int dx^{\mu }\ \frac{\partial \eta _{\rho }}{\partial \alpha }\left\{ 
\frac{\partial \mathcal{L}}{\partial \eta _{\rho }}-\partial _{\nu }\left( 
\frac{\partial \mathcal{L}}{\partial \left( \partial _{\nu }\eta _{\rho
}\right) }\right) \right\} +\int dx^{\mu }\partial _{\nu }\left( \frac{%
\partial \mathcal{L}}{\partial \left( \partial _{\nu }\eta _{\rho }\right) }%
\frac{\partial \eta _{\rho }}{\partial \alpha }\right) ~~~~~.  \label{2.20}
\end{eqnarray}
\end{solution}
